\thispagestyle{empty}
\pdfbookmark[0]{Edition and sources}{edition-and-sources}
\begin{center}
\textsc{Carl Friedrich Gauss}\par\smallskip
{\Large Theory of the Attraction of Homogeneous Elliptic Spheroidal Bodies\par}
\vspace{0pt}{\fontsize{11}{12}\selectfont On a New General Fundamental Law of Mechanics\par}
\vspace{0pt}{\fontsize{11}{12}\selectfont General Principles of the Theory of the Shape of Fluids in Equilibrium\par}
\vspace{0pt}{\fontsize{11}{12}\selectfont The Intensity of the Earth's Magnetic Force Expressed in Absolute Units\par}
\vspace{0pt}{\fontsize{11}{12}\selectfont General Theory of Terrestrial Magnetism\par}
\vspace{0pt}{\fontsize{11}{12}\selectfont General Theorems on Forces of Attraction and Repulsion\par}
\vspace{0pt}{\fontsize{11}{12}\selectfont Dioptrical Investigations\par}
\vspace{0pt}{\fontsize{11}{12}\selectfont Notices of the Author’s Own Papers\par}
\vspace{0pt}{\fontsize{11}{12}\selectfont Introduction to the Magnetic Association's 1836 Observations\par}
\vspace{0pt}{\fontsize{11}{12}\selectfont A New Aid for Magnetic Observations\par}
\vspace{0pt}{\fontsize{11}{12}\selectfont A New Instrument for Measuring Changes in Magnetic Intensity\par}
\vspace{0pt}{\fontsize{11}{12}\selectfont Instructions for Determining the Duration of a Magnetic Needle’s Swing\par}
\vspace{0pt}{\fontsize{11}{12}\selectfont On a Way to Make Deflections Easier to Observe\par}
\vspace{0pt}{\fontsize{11}{12}\selectfont On Determining the Constants of the Bifilar Magnetometer\par}
\vspace{0pt}{\fontsize{11}{12}\selectfont Instructions for Determining the Magnetic Effect Exerted by a Bar Magnet at a Distance\par}
\vspace{0pt}{\fontsize{11}{12}\selectfont On Using the Magnetometer to Determine Absolute Declination\par}
\vspace{0pt}{\fontsize{11}{12}\selectfont Observations of Magnetic Inclination in Göttingen\par}
\vspace{0pt}{\fontsize{11}{12}\selectfont Fundamental Equations for the Motion of Heavy Bodies on the Rotating Earth\par}
\vspace{0pt}{\fontsize{11}{12}\selectfont On Achromatic Doublet Objectives\par}
\vspace{0pt}{\fontsize{11}{12}\selectfont Letter from Gauss to Brandes\par}
\vspace{0pt}{\fontsize{11}{12}\selectfont Adjusting the Knife Edges of a Balance\par}
\vspace{0pt}{\fontsize{11}{12}\selectfont Physical Observations\par}
\vspace{0pt}{\fontsize{11}{12}\selectfont Northern Lights on 7 January 1831\par}
\vspace{0pt}{\fontsize{11}{12}\selectfont Magnetic Observatory in Göttingen (August 1834)\par}
\vspace{0pt}{\fontsize{11}{12}\selectfont Observations of Magnetic Variation in Göttingen and Leipzig\par}
\vspace{0pt}{\fontsize{11}{12}\selectfont Magnetic Observations (March 1835)\par}
\vspace{0pt}{\fontsize{11}{12}\selectfont Magnetic Needle in Copenhagen and Milan\par}
\vspace{0pt}{\fontsize{11}{12}\selectfont Letter from Gauss to Schumacher (1836)\par}
\vspace{0pt}{\fontsize{11}{12}\selectfont Procedure for Scheduled Observations\par}
\vspace{0pt}{\fontsize{11}{12}\selectfont Three Years of Daily Declination Observations\par}
\vspace{0pt}{\fontsize{11}{12}\selectfont Plots and Numerical Readings for the Scheduled Observations of 1836\par}
\vspace{0pt}{\fontsize{11}{12}\selectfont Plots and Numerical Readings for the Scheduled Observations of 1837\par}
\vspace{0pt}{\fontsize{11}{12}\selectfont On Wilkes’s Observations near the Magnetic South Pole\par}
\vspace{0pt}{\fontsize{11}{12}\selectfont Review of Benzenberg’s On Dalton’s Theory\par}
\vspace{0pt}{\fontsize{11}{12}\selectfont Review of Fischer’s Instructions for Making Artificial Magnets\par}
\vspace{0pt}{\fontsize{11}{12}\selectfont Review of the Magnetic Association’s 1836 Results\par}
\vspace{0pt}{\fontsize{11}{12}\selectfont The Mathematical and Physical Institute in Marburg\par}
\vspace{0pt}{\fontsize{11}{12}\selectfont On the Mathematical Theory of Electrodynamic Effects, §§ 1–21, Spherical Harmonics, On Using the Comparator, Effect of a Luminous Point, Schering’s Remarks, Historical Contents and Imprint\par}
\smallskip
AI-translated from Latin and German, with AI-produced typesetting.\\English edition
\end{center}
\clearpage
\begingroup\small
\medskip
\noindent The complete treatise includes all fourteen articles, the historical note and the addendum. The volume title pages are included. The source is the Latin text in Carl Friedrich Gauss, \textit{Werke}, volume V, edited by the Royal Society of Sciences at Göttingen (1867), pp.~3--22, with the original title matter. The title-page vignette is reproduced from the second printing (Göttingen, 1877) digitized by the \href{https://archive.org/details/Werkecarlf05gausrich/page/n8/mode/1up}{University of California Libraries}.
\par\medskip
\noindent \href{https://github.com/KokunoYumeto/gauss-collected-works/releases/tag/gauss-original-language-2026-10-24}{Original-language edition}
\par\medskip
\noindent\textbf{Historical versions and related exposition}\par\smallskip
\noindent The 1813 German version in \textit{Monatliche Correspondenz zur Beförderung der Erd- und Himmels-Kunde}, vol.~28, pp.~37--41, and the German version in Albert Wangerin (ed.), \textit{Ueber die Anziehung homogener Ellipsoide} (Leipzig, 1890), pp.~50--74, were consulted. Arthur Cayley’s \textit{On Gauss’s Method for the Attraction of Ellipsoids} (1857), as reprinted in \textit{Collected Mathematical Papers}, vol.~III (1890), pp.~25--28, is a separate exposition.
\par\medskip
\noindent \href{https://www.digitale-sammlungen.de/de/view/bsb10538617?page=43}{1813}\quad\href{https://archive.org/details/ueberdieanziehu00wanggoog}{1890}\quad\href{https://rcin.org.pl/dlibra/publication/184416/edition/170906/on-gauss-method-for-the-attraction-of-ellipsoids-cayley-arthur-1821-1895?language=en}{Cayley}
\par\medskip
\noindent\textbf{Mechanics}\par\smallskip
\noindent The complete 1829 paper, including both of Gauss’s notes, is AI-translated from German. It appeared in \textit{Journal für die reine und angewandte Mathematik}, vol.~4 (1829), pp.~232--235, and is included here from \textit{Werke}, volume V, pp.~25--28, with its title page. For comparison, D. H. Delphenich’s English translation is available as a standalone PDF and in his \textit{Selected Papers on Gauss’s Principle of Least Constraint} (Neo-Classical Press, 2020).
\par\medskip
\noindent \href{https://doi.org/10.1515/crll.1829.4.232}{1829 original}\quad\href{https://www.neo-classical-physics.info/uploads/3/4/3/6/34363841/gauss_-_princ_least_constraint.pdf}{Delphenich}\quad\href{https://www.neo-classical-physics.info/uploads/3/4/3/6/34363841/selecta_on_least_constraint.pdf#page=95}{2020 collection}
\par\medskip
\Needspace{4\baselineskip}
\noindent\textbf{Fluid equilibrium}\par\smallskip
\noindent The complete Latin treatise is AI-translated with all 34 articles and Gauss’s four notes, from \textit{Werke}, volume V, pp.~31--77, with its 1830 title page. Rudolf Heinrich Weber’s 1903 German translation, edited by Heinrich Weber in Ostwald’s series (no.~135), is a separate historical version.
\par\smallskip
\noindent \href{https://archive.org/details/allgemeinegrund00gausgoog}{Weber’s German edition (1903)}
\par\endgroup
\clearpage
